\section{Why MLOps?}

% --- SLIDE 1: THE FRAMING ---
\begin{frame}{Everything So Far Ended at \texttt{model.fit()}}

    \centering
    \large
    \textit{Every deck in this course stops at the same place: a trained model, in a notebook, on your laptop.}

    \vspace{0.36em}

    \begin{columns}[T]
        \begin{column}{0.46\textwidth}
            \begin{tcolorbox}[colback=blue!5, colframe=blue!50, title=What we have taught you]
                \scriptsize
                \begin{itemize}
                    \item Split the data, pick a metric
                    \item Fit a model, tune it a bit
                    \item Report a validation score
                \end{itemize}
            \end{tcolorbox}
        \end{column}
        \begin{column}{0.46\textwidth}
            \begin{tcolorbox}[colback=red!5, colframe=red!50, title=What nobody taught you]
                \scriptsize
                \begin{itemize}
                    \item Which run produced that number?
                    \item Can you rebuild it in six months?
                    \item Who calls the model, and what happens when it is wrong?
                    \item How will you know it stopped working?
                \end{itemize}
            \end{tcolorbox}
        \end{column}
    \end{columns}

    \vspace{0.3em}
    \scriptsize
    \textbf{MLOps} is the set of practices that answers the right-hand column.

\end{frame}
